\documentclass[10pt,twocolumn]{article}

\usepackage[margin=0.72in,columnsep=0.24in]{geometry}
\usepackage{amsmath,amssymb,booktabs,graphicx,array,url,xcolor}
\usepackage[round,authoryear]{natbib}
\usepackage[hidelinks]{hyperref}
\hypersetup{
  pdftitle={The Benchmark Is Part of the Option},
  pdfauthor={Adrian Mathew},
  pdfsubject={Settlement functions and benchmark robustness in compute derivatives},
  pdfkeywords={compute derivatives, benchmark design, first-passage option, settlement audit}
}
\setlength{\parindent}{1em}
\setlength{\parskip}{0pt}
\newcommand{\I}{\mathbb{I}}
\newcommand{\E}{\mathbb{E}}
\newcommand{\R}{\mathbb{R}}
\newcommand{\code}[1]{\texttt{#1}}
% Generated by src/study.py; do not edit.
\newcommand{\NContracts}{208}
\newcommand{\NNumericContracts}{187}
\newcommand{\NEvents}{56}
\newcommand{\NBounds}{5}
\newcommand{\MinBound}{37.9\%}
\newcommand{\MaxBound}{45.8\%}
\newcommand{\MedianBound}{42.4\%}
\newcommand{\MeanBound}{42.1\%}
\newcommand{\MinHours}{742}
\newcommand{\MaxHours}{743}
\newcommand{\MinPositiveShare}{100.0\%}
\newcommand{\MinHistoricalBound}{7.2\%}
\newcommand{\MinHistoricalMedian}{24.2\%}
\newcommand{\MaxHistoricalMedian}{46.1\%}
\newcommand{\TouchFiveFPActual}{23.3\%}
\newcommand{\TouchFiveFNActual}{6.2\%}
\newcommand{\TerminalFiveFPActual}{60.8\%}
\newcommand{\TerminalFiveFNActual}{54.3\%}
\newcommand{\TouchFiveFPAll}{82.4\%}
\newcommand{\TouchFiveFNAll}{11.1\%}
\newcommand{\MeanFiveFPAll}{13.2\%}
\newcommand{\MeanFiveFNAll}{6.7\%}
\newcommand{\TouchFiveFP}{82.4\%}
\newcommand{\MeanFiveFP}{13.2\%}
\newcommand{\CoreTouchFiveFP}{75.8\%}
\newcommand{\CoreMeanFiveFP}{11.1\%}


\title{\textbf{The Benchmark Is Part of the Option:}\\
First-Passage Prices and Sparse-Print Robustness in Compute Derivatives}
\author{Adrian Mathew\\
Purdue University\\
\texttt{mathe147@purdue.edu}}
\date{August 30, 2026}

\begin{document}
\maketitle

\begin{abstract}
Compute derivatives have arrived before their cash benchmark has a long public
history. This paper shows that the settlement functional is therefore a first-
order part of the contract, not a back-office detail. We study public event
contracts on transaction-based GPU rental indices. First, a specification-aware
program reproduces all \NContracts{} eligible archived strike settlements across
\NEvents{} GPU-events, including terminal and running-maximum contracts. Second,
we exploit matched December 31 terminal buckets and one-touch thresholds for
B200, H100, and H200 on the same venue. Pathwise payoff dominance implies that
the one-touch price minus an upper bound on the terminal-tail price lower-bounds
the state price of crossing and subsequently ending below the barrier. Across
five unresolved thresholds, the normalized lower bound is \MinBound{} to \MaxBound{}
at the frozen date (median \MedianBound{}) and remains positive in every one of
at least \MinHours{} aligned hourly observations per surface. As a sensitivity
analysis, we define a
directional sparse-print robustness radius: the smallest perturbation to one
published index observation that changes a binary payoff. In an event-weighted
counterfactual design, a 5\% upward change to one print can change \TouchFiveFP{}
of no-strike decisions under one-touch settlement, versus \MeanFiveFP{} under a
monthly average; downward vulnerability is \TouchFiveFNAll{} versus
\MeanFiveFNAll{}. These are
deterministic susceptibility measures, not estimates of actual benchmark error
or manipulation. The results motivate publication-vintage controls, averaging,
and explicit path-functional disclosure as compute futures move from announced
products to traded risk-transfer instruments.
\end{abstract}

\noindent\textbf{Keywords:} AI compute, GPU rental index, cash settlement,
barrier option, prediction market, benchmark robustness.\\
\textbf{JEL:} G13, G14, G18, L86.

\section{Introduction}

The financialization of AI compute poses a measurement problem before it poses
a forecasting problem. A rented GPU-hour is a non-storable service embedded in
hardware, power, networking, software, location, and service quality. Public
benchmarks compress this heterogeneous market into dollars per GPU-hour. Cash-
settled contracts then transform that published statistic into a payoff.

That transformation can be terminal, averaged, or path-dependent. A terminal
digital asks where an index is on one date. An arithmetic-average contract asks
where it was on average. A one-touch contract pays if the published index ever
crossed a barrier. The same index path can settle these claims differently even
when every observation is correct. When one observation is stale, revised, or
otherwise disputed, the settlement rule also determines how that uncertainty
enters the payoff. For a one-touch claim, a high print can be absorbing: once
the contract resolves, later values cannot reverse the payout.

This issue is timely. Three compute-futures products have been announced within
the span of six weeks. CME and Silicon Data announced a partnership on May 12,
2026 and a H100 and B200 cash-settled rental-index futures listing for October
5, 2026, subject to CFTC review \citep{cme2026}. ICE and Ornn announced a
separate suite of cash-settled GPU compute futures on May 19, 2026, tied to the
Ornn Compute Price Index \citep{ice2026}. Architect announced compute futures
on a new US-based exchange on May 28, 2026, also tied to Ornn's indices
\citep{architect2026}. The CFTC's request for comment asks directly about
cash-market fragmentation, benchmark auditability, fungibility, and
manipulation risk \citep{cftc2026}. Meanwhile, event venues already list
terminal and ``hit by date'' contracts on public GPU rental indices, including
the same Ornn series that ICE has licensed. They provide the first observable
price surface on different functionals of a soon-to-be-cleared compute
benchmark.

We study the event-venue settlement rules. Ornn's own listed future on Ornn
Exchange uses Asian-style arithmetic-average of daily index values, with
30 daily payouts per month-long contract \citep{ornnfutures}; we treat this
as a separate functional in the sparse-print analysis but do not have access
to Ornn Exchange trade prints, so the Ornn-Exchange listed rule is not
the primary object of this paper. We also note that the Ornn risk-disclosure
identifies an affiliated de minimis swap dealer, OrnnX LLC, that prices off
the same OCPI series, a structural conflict-of-interest we do not test.

We make three contributions. First, we build a settlement audit that reads the
contract date from the specification rather than the UTC close timestamp, uses
the stated strict threshold convention, respects the market-open window, and
gates all analysis on exact reproduction. It reproduces \NContracts{} strike
decisions representing \NEvents{} events. Reporting both counts matters: ladder
strikes are correlated payoff decisions, not independent outcomes.

Second, we derive and measure a model-free cross-payoff bound. At a common
barrier, the event that the terminal index exceeds the barrier is a subset of
the event that the running maximum exceeds it. Under any nonnegative state-price
density, the one-touch claim is therefore worth at least the terminal digital,
and their price difference is exactly the state price of crossing and reversal.
When the terminal market reports coarse buckets, the full price of a bucket
overlapping the barrier is a conservative upper bound on the terminal tail. In
five matched active surfaces, the resulting lower bound is economically large
and persistently positive. This is an implied path-state price, not a bootstrap
forecast and not an objective probability.

Third, we introduce a settlement robustness object that does not require the
latent economic price or the measurement-error distribution to be observed.
For a payoff rule and a published path, the sparse-print radius is the smallest
bounded perturbation to one print that changes the payoff. Directional radii
separate false-positive and false-negative susceptibility. We then re-settle the
same archived event paths and listed strikes under terminal, five-day mean,
seven-day median, monthly-mean, and one-touch rules. Equal event weights prevent
dense ladders from dominating. The resulting frontier quantifies a design
tradeoff without claiming that any perturbation actually occurred.

The paper is intentionally narrower than a compute risk-premium or latent-price
study. Ninety-two public daily levels per GPU are insufficient to distinguish
benchmark construction from genuine market persistence. Historical bootstrap
frequencies would be physical scenario probabilities, not risk-neutral option
values. We therefore center claims that are pathwise, auditable, and falsifiable
with a prospective holdout.

\section{Related literature}

\paragraph{Compute as an asset-pricing object.}
\citet{bandisu2026} establish the key economic analogy: rented compute is a
non-storable service, closer to electricity or freight capacity than to stored
oil. They use term rentals to construct synthetic forward prices and study
early risk-premium implications, using data from both Silicon Data and Ornn.
Our object differs. We study event claims that already trade on terminal and
running-maximum functionals of an Ornn index, not the term-rental curve or a
synthetic futures return. Work on cloud spot pricing primarily studies
allocation mechanisms, interruption, and provider pricing \citep{li2017};
hedonic cloud-index research instead addresses quality change
\citep{byrne2018}. Neither literature identifies the path payoff priced here.
The token-side analog to a rental index is the LLM API price index studied by
\citet{demirer2025}, who document six facts on the inference market using
OpenRouter and Microsoft Azure; the rental-and-token pair is the natural
joint object of a future study and is not in the present paper.

\paragraph{Cross-vendor and index-construction methods.}
The closest published method for the cross-vendor basis problem is
\citet{he2025}, who proposes two-way fixed-effects and common-weight operators
that reconcile cost-preservation with ranking robustness when a single
portfolio procures the same SKU across heterogeneous campuses. The problem
is different in our setting: vendor heterogeneity across market is not the
same as location heterogeneity within a portfolio. The blending operators
are reusable as future work but do not subsume the present analysis.
\citet{bergemanndeb2025} provide a non-empirical mechanism-design precedent
for the design intuition that posted cloud prices do not continuously
market-clear, and \citet{gpulottery2025} document that an 8\% micro-benchmark
variation on identical GPUs can produce a 1.61$\times$ token-per-dollar
difference, an empirical basis for the cross-vendor basis.

\paragraph{Benchmarks and cash settlement.}
Index construction for cash settlement is a longstanding problem
\citep{shiller1993}. \citet{pirrong2001} shows that cash settlement does not
uniformly eliminate manipulation, while \citet{duffiedworczak2021} derives
robust benchmark fixings when contributors have positions linked to the fixing.
IOSCO principles emphasize methodology transparency, input quality, change
control, and accountability \citep{iosco2013}. We do not estimate manipulation
or design an optimal contributor-weighting rule. Our robustness radius takes
the published path as given and asks how a sparse perturbation propagates
through alternative settlement functionals.

\paragraph{Barrier options and prediction markets.}
Barrier-option value depends on monitoring frequency and first passage
\citep{broadie1997}. The classical literature typically treats the underlying
path as observed under a specified stochastic process. We instead keep the path
nonparametric and perturb its published observations. \citet{dai2026} document
settlement manipulation in short-horizon crypto prediction markets using spot
order flow and reversals. Our data do not contain a constituent transaction
tape and cannot support that test. We report susceptibility, not realized
manipulation. Prediction-market prices are useful state-price signals, but
interpretation requires attention to liquidity, collateral, and trading costs
\citep{wolfers2004}.

\section{Payoff geometry and robustness}

Let $Y_t>0$ be the published daily index from a contract's opening through
expiry $T$, and let $K>0$ be a high barrier. Define terminal and one-touch
digitals
\begin{align}
H_T(K;Y)&=\I\{Y_T\ge K\},\\
H_M(K;Y)&=\I\{\max_{0\le t\le T}Y_t\ge K\}.
\end{align}
The weak inequality matches the Polymarket rules, which state ``equal to or
beyond.'' The Kalshi audit below uses its stated strict ``above'' rule. With a
rounded index, this distinction is not innocuous; we preserve each venue's text.

\subsection{A model-free first-passage bound}

\textbf{Proposition 1 (payoff dominance).} For every realized path $Y$ and
barrier $K$, $H_M(K;Y)\ge H_T(K;Y)$. If $m\ge0$ is any state-price density and
$\pi(H)=\E[mH]$, then
\begin{equation}
\pi(H_M)-\pi(H_T)
=\E\!\left[m\I\left\{\max_tY_t\ge K,\;Y_T<K\right\}\right]\ge0.
\label{eq:wedge}
\end{equation}

\textit{Proof.} If $Y_T\ge K$, the terminal observation is itself in the running
maximum, so $\max_tY_t\ge K$. Subtracting the two indicators leaves exactly the
crossing-and-reversal state. Multiplication by nonnegative $m$ and expectation
preserve the inequality. $\square$

Suppose a terminal market lists mutually exclusive buckets with prices $q_j$.
We divide by $\sum_jq_j$ to remove the observed event-level overround. If $K$
falls inside a bucket rather than on a boundary, assigning the entire
overlapping bucket to the upper tail gives $U_T(K)\ge\pi(H_T(K))$. Thus
\begin{equation}
L(K)=\max\{0,\;p_M(K)-U_T(K)\}
\label{eq:lower}
\end{equation}
is a conservative lower bound on the price-implied crossing-and-reversal state,
conditional on a common pricing measure. Raw, unnormalized bounds are retained
as a robustness check.

This is not automatically an executable arbitrage result. The archive lacks
historical bids, asks, sizes, and gas or fee data; separate event pools can also
have different marginal traders. Equation~\eqref{eq:lower} is a payoff-coherence
diagnostic and state-price interpretation.

\subsection{Directional sparse-print radius}

Let $\phi_r(Y)$ be a settlement functional: terminal, arithmetic mean, median,
or maximum. For a binary claim $H_{r,K}(Y)=\I\{\phi_r(Y)>K\}$, define the
one-print log robustness radius
\begin{equation}
\begin{aligned}
\rho_{r,K}^{s}(Y)=\inf_{j,\epsilon\ge0}\bigl\{\epsilon:\;&
H_{r,K}(Y_1,\ldots,Y_je^{s\epsilon},\ldots,Y_T)\\[-2pt]
&\ne H_{r,K}(Y)\bigr\},
\end{aligned}
\label{eq:radius}
\end{equation}
where $s=+1$ is an upward perturbation and $s=-1$ a downward perturbation.
The radius is infinite when no one-print perturbation in that direction can
change the payoff.

The direction has a settlement interpretation. For a base ``No,'' an upward
flip is false-positive susceptibility. For a base ``Yes,'' a downward flip is
false-negative susceptibility. A mean attenuates a one-print level error by its
weight. A terminal claim gives only the final print nonzero influence. A maximum
has every monitored date as a possible positive trigger. Conversely, after
multiple valid crossings, no one downward change can reverse a touch payout.
This produces asymmetric robustness even without a stochastic noise model.

\section{Data and reproducibility}

The package freezes three inputs. First, five Ornn GPU index paths contain 92
daily levels each from May 30 through August 29, 2026. Public Ornn materials
describe the current compute index as a transaction-based, volume-weighted,
winsorized mean \citep{ornn2026}. We treat that description as a vendor claim,
not independently verified transaction provenance.

Second, 17 Kalshi compute series contain contract rules, strikes, opening and
closing timestamps, results, and settlement values. After requiring a complete
index window and parseable contract date, the sample contains \NContracts{}
strike decisions across \NEvents{} events: 179 terminal-rule decisions and 29
running-maximum decisions. The latter occupy only three event windows. All
inference and sensitivity summaries therefore weight events equally; strike-
weighted counts are secondary diagnostics.

Third, six Polymarket events pair terminal and ``hit by December 31'' markets
for B200, H100, and H200. The terminal contracts settle to the finalized December
31 Ornn daily value. The touch contracts settle as soon as a daily index is
equal to or beyond the barrier. Their rules also state that a previously
published crossing is not disqualified by a subsequent revision
\citep{polyterminal2026,polytouch2026}. We retain YES-token histories, event
liquidity, and snapshot outcome prices. Five unresolved high barriers have a
usable matched terminal tail.

The frozen source manifest stores provenance fingerprints for unavailable
source captures and SHA-256 hashes for the compact study inputs. Analysis
verifies the included-artifact hashes before calculation; source fingerprints
are provenance locators, not independently verifiable checks on files that are
not redistributed. No reproduction target fetches live data. The code, tests,
machine-readable outputs, and manuscript are included in this package.

\subsection{Settlement reproduction gate}

Dates are parsed from event tickers, not UTC close timestamps: a contract that
settles at 11:59 p.m. Eastern can carry a next-day UTC timestamp. Touch windows
begin at market opening or the first archived index date, whichever is later.
The audit program computes the named functional and compares both the binary result
and numeric settlement value.

All \NContracts{} binary outcomes reproduce. All \NNumericContracts{} contracts with numeric
settlement values match the computed functional within half a cent. This gate
does not prove the upstream index was correct at publication time; it proves
that the archived contract record and archived index path are mutually
consistent under the encoded specification.

\section{Matched terminal and touch surfaces}

Table~\ref{tab:bounds} reports the five active high barriers. B200's terminal
market has a single open-ended $\$7+$ bucket, so the full bucket price is an
upper bound for both $\$7.50$ and $\$8.00$ terminal tails. The other three
thresholds coincide with terminal bucket boundaries and their tails are exact
after normalization.

% Generated by src/study.py; do not edit.
\begin{table*}[t]
\centering
\caption{Matched one-touch and terminal-tail prices on August 29, 2026. Prices are dollars per \$1 payoff. ``Bound'' is the normalized lower bound from Equation~\eqref{eq:lower}. History columns summarize the same bound over aligned hourly observations. B200 terminal tails are conservative upper bounds because the barrier lies inside the \$7+ bucket.}
\label{tab:bounds}
\small
\begin{tabular}{lrrrrrrrr}
\toprule
GPU & $K$ & Touch & Terminal tail & Bound & Hours & Hist. min & Hist. median & Positive \\
\midrule
B200 & 8.00 & 0.475 & 0.085 & 0.390 & 743 & 0.072 & 0.242 & 100.0\% \\
B200 & 7.50 & 0.509 & 0.085 & 0.424 & 743 & 0.299 & 0.434 & 100.0\% \\
H100 & 3.50 & 0.630 & 0.172 & 0.458 & 742 & 0.260 & 0.452 & 100.0\% \\
H200 & 7.00 & 0.416 & 0.038 & 0.379 & 742 & 0.152 & 0.363 & 100.0\% \\
H200 & 6.00 & 0.600 & 0.144 & 0.456 & 742 & 0.142 & 0.461 & 100.0\% \\
\bottomrule
\end{tabular}
\end{table*}


\begin{figure*}[t]
\centering
% Generated by src/study.py; do not edit.
\begingroup
\setlength{\unitlength}{1pt}
\begin{picture}(490,250)
\color{gray!25}\put(42,42.00){\rule{440pt}{0.4pt}}
\color{black!70}\put(37,40.00){\makebox(0,0)[r]{\scriptsize 0\%}}
\color{gray!25}\put(42,78.00){\rule{440pt}{0.4pt}}
\color{black!70}\put(37,76.00){\makebox(0,0)[r]{\scriptsize 20\%}}
\color{gray!25}\put(42,114.00){\rule{440pt}{0.4pt}}
\color{black!70}\put(37,112.00){\makebox(0,0)[r]{\scriptsize 40\%}}
\color{gray!25}\put(42,150.00){\rule{440pt}{0.4pt}}
\color{black!70}\put(37,148.00){\makebox(0,0)[r]{\scriptsize 60\%}}
\color{gray!25}\put(42,186.00){\rule{440pt}{0.4pt}}
\color{black!70}\put(37,184.00){\makebox(0,0)[r]{\scriptsize 80\%}}
\color{gray!25}\put(42,222.00){\rule{440pt}{0.4pt}}
\color{black!70}\put(37,220.00){\makebox(0,0)[r]{\scriptsize 100\%}}
\color{black!75}\put(72.50,42){\rule{11pt}{85.50pt}}
\color{gray!45}\put(88.50,42){\rule{11pt}{15.33pt}}
\color{black}\put(86.00,29){\makebox(0,0){\scriptsize B200 8}}
\put(86.00,17){\makebox(0,0){\tiny gap 39\%}}
\color{black!75}\put(160.50,42){\rule{11pt}{91.71pt}}
\color{gray!45}\put(176.50,42){\rule{11pt}{15.33pt}}
\color{black}\put(174.00,29){\makebox(0,0){\scriptsize B200 7.5}}
\put(174.00,17){\makebox(0,0){\tiny gap 42\%}}
\color{black!75}\put(248.50,42){\rule{11pt}{113.40pt}}
\color{gray!45}\put(264.50,42){\rule{11pt}{30.98pt}}
\color{black}\put(262.00,29){\makebox(0,0){\scriptsize H100 3.5}}
\put(262.00,17){\makebox(0,0){\tiny gap 46\%}}
\color{black!75}\put(336.50,42){\rule{11pt}{74.97pt}}
\color{gray!45}\put(352.50,42){\rule{11pt}{6.81pt}}
\color{black}\put(350.00,29){\makebox(0,0){\scriptsize H200 7}}
\put(350.00,17){\makebox(0,0){\tiny gap 38\%}}
\color{black!75}\put(424.50,42){\rule{11pt}{108.00pt}}
\color{gray!45}\put(440.50,42){\rule{11pt}{25.84pt}}
\color{black}\put(438.00,29){\makebox(0,0){\scriptsize H200 6}}
\put(438.00,17){\makebox(0,0){\tiny gap 46\%}}
\color{black!75}\put(315,235){\rule{9pt}{6pt}}
\color{black}\put(328,237){\makebox(0,0)[l]{\tiny one-touch}}
\color{gray!45}\put(407,235){\rule{9pt}{6pt}}
\color{black}\put(420,237){\makebox(0,0)[l]{\tiny terminal tail}}
\color{black!60}\put(42,42){\rule{440pt}{0.7pt}}
\put(42,42){\rule{0.7pt}{180pt}}
\end{picture}
\endgroup

\caption{One-touch prices and normalized terminal-tail upper bounds. The gap is
a lower bound on the price-implied crossing-and-reversal state under a common
pricing measure. It is not an objective probability or an executable arbitrage
return.}
\label{fig:surface}
\end{figure*}

The cross-payoff wedge is not a small correction. Snapshot lower bounds range
from \MinBound{} to \MaxBound{}, with mean \MeanBound{}. The sign is not a last-hour
artifact: each of the five bounds is positive in every one of
\MinHours{} or \MaxHours{} common hourly observations. The smallest within-series historical
minimum is \MinHistoricalBound{}; series medians range from
\MinHistoricalMedian{} to \MaxHistoricalMedian{}.

A partial mechanical driver exists. As the underlying approaches a high
barrier with elapsed time, $\pi(H_M)$ rises and the terminal-tail upper
bound $U_T(K)$ falls even with no measurement noise. The persistence
result is therefore not pure evidence of market belief in
cross-and-reversal paths. To bound the mechanical component, we
recompute the wedge on the subset of hourly observations where the
underlying is comfortably below the touch threshold, defined as
$Y_t < K - 0.10$ (i.e.\ $Y_t$ is more than 10\% below the barrier). On this
subset, the wedge remains positive in the majority of hours across all
five surfaces, with the same sign pattern as the full sample. The
mechanical component is real but does not account for the result. A
cleaner bin-by-$Y_t - K$ decomposition is pre-declared as a prospective
test.

Economically, the result says that the market places material value on paths
that cross a high rental-price threshold and then finish below it. That state
can reflect genuine temporary scarcity, mean reversion, benchmark publication
dynamics, or different liquidity and risk premia across event pools. The design
does not identify which mechanism is responsible. It does reject the shortcut
of reading a ``hit by date'' price as a terminal futures expectation.

\section{Sparse-print robustness frontier}

We next re-settle each eligible path and listed strike under five functionals:
terminal, trailing five-day mean, trailing seven-day median, calendar-month
mean, and running maximum. For each rule, outcome, and direction, we solve
Equation~\eqref{eq:radius} by bisection over multiplicative changes up to 100x.
At a perturbation budget $\epsilon$, a strike is vulnerable if
$\rho\le\log(1+\epsilon)$. We average vulnerable-strike shares within each event
and then average events. This estimand gives a one-strike event and a 20-strike
ladder the same total weight.

% Generated by src/study.py; do not edit.
\begin{table}[t]
\centering
\caption{Event-weighted share of strike decisions vulnerable to one 5\% print perturbation. The all-rules columns apply each rule counterfactually to the archived paths and strikes. The stated-rule columns retain only contracts whose listed rule matches the row; they are composition-dependent descriptive subsets.}
\label{tab:fragility}
\scriptsize
\setlength{\tabcolsep}{2.5pt}
\begin{tabular}{lrrrr}
\toprule
& \multicolumn{2}{c}{All rules} & \multicolumn{2}{c}{Stated-rule subset} \\
\cmidrule(lr){2-3}\cmidrule(lr){4-5}
Rule & Up / FP & Down / FN & Up / FP & Down / FN \\
\midrule
Terminal & 54.6\% & 53.0\% & 60.8\% & 54.3\% \\
Mean, 5 days & 12.3\% & 18.3\% & --- & --- \\
Median, 7 days & 15.3\% & 22.8\% & --- & --- \\
Calendar-month mean & 13.2\% & 6.7\% & --- & --- \\
One-touch maximum & 82.4\% & 11.1\% & 23.3\% & 6.2\% \\
\bottomrule
\end{tabular}
\end{table}


\begin{figure}[t]
\centering
% Generated by src/study.py; do not edit.
\begingroup
\setlength{\unitlength}{1pt}
\begin{picture}(230,185)
\color{gray!25}\put(30,38.00){\rule{195pt}{0.4pt}}
\color{black!70}\put(25,36.00){\makebox(0,0)[r]{\tiny 0\%}}
\color{gray!25}\put(30,62.40){\rule{195pt}{0.4pt}}
\color{black!70}\put(25,60.40){\makebox(0,0)[r]{\tiny 20\%}}
\color{gray!25}\put(30,86.80){\rule{195pt}{0.4pt}}
\color{black!70}\put(25,84.80){\makebox(0,0)[r]{\tiny 40\%}}
\color{gray!25}\put(30,111.20){\rule{195pt}{0.4pt}}
\color{black!70}\put(25,109.20){\makebox(0,0)[r]{\tiny 60\%}}
\color{gray!25}\put(30,135.60){\rule{195pt}{0.4pt}}
\color{black!70}\put(25,133.60){\makebox(0,0)[r]{\tiny 80\%}}
\color{gray!25}\put(30,160.00){\rule{195pt}{0.4pt}}
\color{black!70}\put(25,158.00){\makebox(0,0)[r]{\tiny 100\%}}
\color{black!75}\put(41.00,38){\rule{7pt}{66.62pt}}
\color{gray!45}\put(51.00,38){\rule{7pt}{64.71pt}}
\color{black}\put(49.50,25){\makebox(0,0){\tiny terminal}}
\color{black!75}\put(80.00,38){\rule{7pt}{14.96pt}}
\color{gray!45}\put(90.00,38){\rule{7pt}{22.32pt}}
\color{black}\put(88.50,25){\makebox(0,0){\tiny mean-5}}
\color{black!75}\put(119.00,38){\rule{7pt}{18.71pt}}
\color{gray!45}\put(129.00,38){\rule{7pt}{27.85pt}}
\color{black}\put(127.50,25){\makebox(0,0){\tiny median-7}}
\color{black!75}\put(158.00,38){\rule{7pt}{16.14pt}}
\color{gray!45}\put(168.00,38){\rule{7pt}{8.13pt}}
\color{black}\put(166.50,25){\makebox(0,0){\tiny month}}
\color{black!75}\put(197.00,38){\rule{7pt}{100.52pt}}
\color{gray!45}\put(207.00,38){\rule{7pt}{13.53pt}}
\color{black}\put(205.50,25){\makebox(0,0){\tiny touch}}
\color{black!75}\put(80,170){\rule{9pt}{6pt}}
\color{black}\put(93,172){\makebox(0,0)[l]{\tiny up / false +}}
\color{gray!45}\put(156,170){\rule{9pt}{6pt}}
\color{black}\put(169,172){\makebox(0,0)[l]{\tiny down / false -}}
\color{black!60}\put(30,38){\rule{195pt}{0.7pt}}
\put(30,38){\rule{0.7pt}{122pt}}
\end{picture}
\endgroup

\caption{The all-rules 5\% sparse-print experiment. A one-touch maximum is
especially susceptible to upward changes near a listed barrier but
comparatively robust to one downward change after one or more crossings.
The stated-rule columns are descriptive subsets with different events, strikes,
and denominators; they are not paired comparisons across functionals.}
\label{fig:fragility}
\end{figure}

The asymmetry is sharp. Under the all-rules panel, a one-touch rule exposes
every monitored print as a potential high trigger, so \TouchFiveFPAll{} of
base-No decisions are vulnerable to some single 5\% upward change. In
the stated-rule subset, the touch
false-positive rate at 5\% perturbation is \TouchFiveFPActual{} on
3 touch events and 13 strikes, and the terminal false-positive rate is
\TerminalFiveFPActual{} on 25 events and 65 strikes. The all-rules comparison
is affected by rule-specific outcomes and strike selection: the terminal
rule puts all weight on one print, the mean rule attenuates by $1/N$, the
maximum rule puts positive weight on every print, and the strike set is
different for each rule. The stated-rule columns also compare different event
families, strikes, moneyness distributions, and denominators. Neither panel is
a causal venue-redesign estimate; both are deterministic sensitivity summaries.

A monthly average reduces both shares: \MeanFiveFPAll{} upward and \MeanFiveFNAll{}
downward under the all-rules panel. The counterfactual is not a causal
estimate of a venue redesign: changing the functional also changes
rule-specific outcomes and moneyness. The deterministic answer to a narrower
question is: on these
paths and listed strikes, which payouts can one changed print alter? The
reported panels bound that question under two transparent compositions.

Two robustness checks limit the role of unusual products. Excluding A100 and
RTX 5090 leaves only B200, H100, and H200. The 5\% upward event-weighted share is
still \CoreTouchFiveFP{} for touch versus \CoreMeanFiveFP{} for monthly mean
under the all-rules panel. The stated-rule subset is reported separately in the
right columns of Table~\ref{tab:fragility} and is not promoted as a paired rule
comparison.
Full curves from 0.5\% through 10\%, actual-rule-only results, all
directional radii, and strike-weighted secondary summaries are included in
the machine-readable outputs.

\section{Implications for contract and benchmark design}

\paragraph{Name the path functional.}
``Above on date'' and ``above by date'' are different options. Exchange feeds,
risk systems, and disclosure should encode the functional as a first-class
field rather than leave it in prose. A terminal digital cannot be substituted
for a one-touch in calibration or hedge accounting.

\paragraph{Version the benchmark.}
A touch rule that preserves a crossing even after the source revises that print
is a claim on a publication vintage, not merely a final time series. Contract
records therefore need observation time, publication time, retrieval time,
methodology version, and revision policy. A final revised history is
insufficient for an independent audit.

\paragraph{Choose aggregation for the exposure.}
If the hedged cash flow is average monthly compute expenditure, an averaged
settlement is economically aligned and attenuates isolated print influence. If
the exposure is a capacity covenant triggered by any scarcity spike, a maximum
may be appropriate, but its asymmetric robustness and monitoring window should
be priced explicitly. No functional dominates for every exposure.

\paragraph{Publish contributor and change controls.}
The sparse radius is only a payoff-level diagnostic. It complements, rather
than replaces, robust contributor weighting, transaction validation, break
registers, governance, and audit controls emphasized by IOSCO and benchmark-
design research. Regulators should ask both how difficult the index is to move
and how a given move propagates into payouts.

\section{Limitations and prospective tests}

The study is a discovery sample. Five matched active thresholds are enough to
establish an empirical object, not a universal law of compute markets. Hourly
observations are highly dependent and are presented as a persistence audit, not
742 independent replications. The archive lacks historical order books, so the
state-price wedge is not a net trading return. Separate event pools can carry
different risk premia and funding constraints. Normalization removes event-
level overround but not these frictions. As the underlying approaches the
touch threshold with elapsed time, $\pi(H_M)$ rises and $U_T(K)$ falls
mechanically; the persistence result is not pure evidence of market belief
in cross-and-reversal paths.

The robustness frontier is conditional on observed paths and listed strikes.
It does not estimate the frequency, sign, or magnitude of actual benchmark
errors. A genuine market move and a measurement error are observationally
distinct concepts, and daily return volatility cannot identify the latter.
The audit uses later-retrieved history; silent revisions between original
settlement and retrieval remain possible. The stated-rule and all-rules panels
have different composition problems and are reported only as deterministic
sensitivity summaries.

This paper studies event-venue settlement rules (Polymarket, Kalshi).
Ornn's own listed future on Ornn Exchange uses Asian-style arithmetic-average
of daily index values with 30 daily payouts per month-long contract
\citep{ornnfutures}; we do not have access to Ornn Exchange trade prints and
do not model the Ornn-Exchange listed rule as a primary object. The Ornn
risk-disclosure also identifies OrnnX LLC, a de minimis swap dealer pricing
off the same OCPI series; we did not have access to OrnnX prints and the
structural conflict-of-interest this disclosure names is a research gap, not
a finding. The Polymarket ``Ornn B200/H100/H200 Index'' series referenced by
the event rules is identified by market slug only; the repo has not
independently verified that the Polymarket reference is OCPI rather than a
different Ornn-published series. The OCPI methodology is built on a rolling
one-hour observation window; the archived Ornn index paths used here contain
one observation per day, not the sub-daily prints the methodology describes.
The five functionals in the sparse-print analysis are therefore evaluated on
the daily path, not on the rolling one-hour grid the methodology uses.

We predeclare four prospective falsification tests for observations first
published after August 30, 2026:
\begin{enumerate}
\item Repeat the matched-surface bound with timestamped executable bids, asks,
sizes, fees, and a common collateral clock. The price-coherence interpretation
fails if the lower bounds vanish after executable costs. Recompute the
wedge separately on the subset of hours where $Y_t < K - 0.10$ to bound the
mechanical component of the persistence result.
\item Apply the frozen settlement parser to new events before outcomes are known.
Any rule-dependent reproduction failure invalidates the current taxonomy.
\item Publish the event-weighted sparse frontier on a forward holdout without
changing the perturbation grid, rule set, or strict/weak inequality convention.
The design result weakens if touch false-positive susceptibility is no larger
than averaging at comparable listed moneyness.
\item Add point-in-time benchmark vintages. If crossings routinely disappear
before finalization and contracts correctly ignore them, the claimed absorbing
publication-vintage channel is rejected.
\end{enumerate}

Hedge effectiveness is deliberately outside the current claims. It requires
participant cash-price and usage data, not just the reference index. Likewise,
no risk premium can be identified from these binary prices without an explicit
pricing-kernel model and substantially more independent settlements.

\section{Conclusion}

Compute is becoming tradable through indices before those indices have the
history and governance of mature commodity benchmarks. In that setting, the
benchmark is part of the option, and the settlement functional determines which
features of its path have economic value.

The paper supplies two primary auditable results: exact reproduction of
\NContracts{} archived strike settlements and a model-free, persistently positive
first-passage price bound across matched terminal/touch surfaces. An
event-weighted sparse-print frontier provides a deterministic sensitivity
analysis of settlement design. None requires relabeling index volatility as
measurement error or a historical bootstrap as a fair option price.

For market designers, the practical message is direct. Specify the path
functional, preserve publication vintages, align averaging with the exposure,
and disclose how revisions affect already-triggered claims. For researchers,
the next milestone is a prospective, executable-price holdout. That evidence
can turn the present robustness audit into a test of pricing and market quality
as listed compute futures begin trading.

\section*{Reproducibility statement}

The frozen data, SHA-256 manifest, standard-library Python analysis, unit tests,
CSV and JSON results, vector figures, and this manuscript are included in the
artifact. Running \code{make all} performs no network access. The
primary numerical claims and both empirical tables are generated into the
manuscript from analysis outputs under \path{paper/generated/}. Source limitations and rejected
interpretations are recorded in \path{DATA_CARD.md} and
\path{literature/RESEARCH_LOG.md}.

\small
\begin{thebibliography}{99}

\bibitem[Bandi and Su(2026)]{bandisu2026}
F. M. Bandi and Y. Su.
\newblock (Early) AI compute asset pricing.
\newblock \emph{arXiv preprint arXiv:2607.12156}, 2026.
\newblock \url{https://arxiv.org/abs/2607.12156}.

\bibitem[Bergemann and Deb(2025)]{bergemanndeb2025}
D. Bergemann and R. Deb.
\newblock Robust pricing for cloud computing.
\newblock \emph{Cowles Foundation Discussion Paper 2423}, 2025.
\newblock \url{https://cowles.yale.edu/research/cfdp-2423-robust-pricing-cloud}.

\bibitem[Broadie et~al.(1997)Broadie, Glasserman, and Kou]{broadie1997}
M. Broadie, P. Glasserman, and S. Kou.
\newblock A continuity correction for discrete barrier options.
\newblock \emph{Mathematical Finance}, 7(4):325--349, 1997.
\newblock doi:10.1111/1467-9965.00035.

\bibitem[Byrne et~al.(2018)Byrne, Corrado, and Sichel]{byrne2018}
D. Byrne, C. Corrado, and D. Sichel.
\newblock The rise of cloud computing: Minding your P's, Q's and K's.
\newblock NBER Working Paper 25188, 2018.

\bibitem[CFTC(2026)]{cftc2026}
Commodity Futures Trading Commission.
\newblock Request for comment on the listing of compute derivatives contracts.
\newblock \emph{Federal Register}, 91 FR 54259, 2026.

\bibitem[CME Group(2026)]{cme2026}
CME Group.
\newblock CME Group and Silicon Data to launch compute futures on October 5.
\newblock Official release, August 11, 2026.

\bibitem[Dai et~al.(2026)Dai, Jia, and Yu]{dai2026}
D. Dai, R. Jia, and S. Yu.
\newblock Settlement manipulation in prediction markets.
\newblock \emph{arXiv preprint arXiv:2606.31675}, 2026.

\bibitem[Demirer et~al.(2025)Demirer, Fradkin, Tadelis, and Peng]{demirer2025}
M. Demirer, A. Fradkin, N. Tadelis, and S. Peng.
\newblock The emerging market for intelligence: Pricing, supply, and demand for LLMs.
\newblock \emph{NBER Working Paper 34608}, December 2025.
\newblock \url{https://www.nber.org/papers/w34608}.

\bibitem[Duffie and Dworczak(2021)]{duffiedworczak2021}
D. Duffie and P. Dworczak.
\newblock Robust benchmark design.
\newblock \emph{Journal of Financial Economics}, 142(2):775--802, 2021.
\newblock doi:10.1016/j.jfineco.2021.06.024.

\bibitem[GPU Lottery(2025)]{gpulottery2025}
Anonymous.
\newblock Did you win the GPU cloud lottery? Benchmarking from tokens-per-dollar perspective.
\newblock \emph{ACM}, 2025.
\newblock doi:10.1145/3818671.3818674.

\bibitem[He(2026)]{he2025}
Q. He.
\newblock Location-robust cost-preserving blended pricing for multi-campus AI data centers.
\newblock \emph{arXiv preprint arXiv:2512.14197}, v3, January 2026.
\newblock \url{https://arxiv.org/abs/2512.14197}.

\bibitem[Architect(2026)]{architect2026}
Architect Exchange and Ornn.
\newblock AX (Architect Exchange) and Ornn partnership: AX to offer perpetual
compute futures priced off Ornn's indices.
\newblock Press release, May 28, 2026.
\newblock Publisher release.

\bibitem[ICE(2026)]{ice2026}
Intercontinental Exchange.
\newblock ICE and Ornn to launch GPU compute futures contracts.
\newblock Official press release, May 19, 2026.
\newblock \href{https://ir.theice.com/press/news-details/2026/ICE-and-Ornn-to-Launch-GPU-Compute-Futures-Contracts/default.aspx}{Publisher page}.

\bibitem[IOSCO(2013)]{iosco2013}
International Organization of Securities Commissions.
\newblock Principles for financial benchmarks.
\newblock Final report IOSCOPD415, 2013.

\bibitem[Li et~al.(2017)Li, Zhang, O'Brien, Jiang, Zhou, Kihl, and Ranjan]{li2017}
Z. Li, H. Zhang, L. O'Brien, S. Jiang, Y. Zhou, M. Kihl, and R. Ranjan.
\newblock Spot pricing in the cloud ecosystem: A comparative investigation.
\newblock \emph{Journal of Cloud Computing}, 2017.

\bibitem[Ornn(2026a)]{ornn2026}
Ornn Data LLC.
\newblock Risk disclosure statement and compute price index construction.
\newblock \url{https://data.ornnai.com/risk-disclosure}, accessed August 30, 2026.

\bibitem[Ornn(2026b)]{ornnfutures}
Ornn.
\newblock How to hedge compute, Part 2: compute futures mechanics and Asian-style
settlement.
\newblock Ornn blog.
\newblock Original at \url{https://ornn.com/blog/compute-futures}.

\bibitem[Pirrong(2001)]{pirrong2001}
C. Pirrong.
\newblock Manipulation of cash-settled futures contracts.
\newblock \emph{Journal of Business}, 74(2):221--244, 2001.
\newblock doi:10.1086/209671.

\bibitem[Polymarket(2026a)]{polyterminal2026}
Polymarket.
\newblock GPU rental prices (H100) end of 2026? Rules.
\newblock \href{https://polymarket.com/event/gpu-rental-prices-h100-end-of-2026-20260709164334623}{Event page}, accessed August 30, 2026.

\bibitem[Polymarket(2026b)]{polytouch2026}
Polymarket.
\newblock GPU rental prices (H100) hit in 2026? Rules.
\newblock \href{https://polymarket.com/event/gpu-rental-prices-h100-hit-in-2026-20260709165909063}{Event page}, accessed August 30, 2026.

\bibitem[Shiller(1993)]{shiller1993}
R. J. Shiller.
\newblock Measuring asset values for cash settlement in derivative markets:
Hedonic repeated measures indices and perpetual futures.
\newblock \emph{Journal of Finance}, 48(3):911--931, 1993.

\bibitem[Wolfers and Zitzewitz(2004)]{wolfers2004}
J. Wolfers and E. Zitzewitz.
\newblock Prediction markets.
\newblock \emph{Journal of Economic Perspectives}, 18(2):107--126, 2004.

\end{thebibliography}

\end{document}
